\chapter{\Lang{Hasil dan Pembahasan}{Results and Discussion}}

\section{\Lang{Hasil Implementasi atau Pelaksanaan Penelitian}{Implementation or Research Execution Results}}
\TemplateTodo{\Lang{Laporkan apa yang benar-benar berhasil dibuat, dikumpulkan, atau dijalankan. Pisahkan hasil faktual dari interpretasi.}{Report what was actually built, collected, or executed. Separate factual results from interpretation.}}

\section{\Lang{Hasil Pengujian}{Evaluation Results}}
\TemplateTodo{\Lang{Sajikan hasil utama sesuai skenario pada Bab III. Tabel harus memiliki judul jelas, satuan, dan keterangan yang cukup.}{Present main results according to the scenarios in Chapter III. Tables should have clear captions, units, and sufficient explanation.}}

\begin{table}[H]
\centering
\caption{\Lang{Contoh tabel hasil pengujian}{Example evaluation result table}}
\begin{tabular}{@{}lcc@{}}
\toprule
\Lang{Skenario}{Scenario} & \Lang{Metrik 1}{Metric 1} & \Lang{Metrik 2}{Metric 2} \\
\midrule
\Lang{Skenario A}{Scenario A} & nilai/value & nilai/value \\
\Lang{Skenario B}{Scenario B} & nilai/value & nilai/value \\
\bottomrule
\end{tabular}
\end{table}

\section{\Lang{Pembahasan}{Discussion}}
\TemplateTodo{\Lang{Jelaskan mengapa hasil tersebut terjadi, hubungannya dengan teori dan penelitian terdahulu, serta implikasinya. Jangan hanya mengulang isi tabel.}{Explain why the results occurred, their relationship to theory and related work, and their implications. Do not merely repeat table contents.}}

\section{\Lang{Keterbatasan Penelitian}{Research Limitations}}
\TemplateTodo{\Lang{Tuliskan keterbatasan secara jujur dan spesifik, misalnya keterbatasan data, waktu, perangkat, metode, validasi, atau cakupan pengguna.}{State limitations honestly and specifically, such as limitations in data, time, hardware, method, validation, or user scope.}}
